\section{分层策略：从你的位置出发}

方法库是通用的，但你的起跑线不一样。这一章把前面的所有东西，按学历光谱的两端分别组装成两套不同的自救方案。

\subsection{学历上层：从「撑住」到「可持续」}

你的问题不是「动不了」，是「靠肾上腺素动，快烧完了」。所以你的自救关键词只有一个：\term{拆弹}——拆掉那台让你撑到现在、也快把你烧穿了的引擎。

\begin{checklist}
  \item \textbf{停止等死线}。死线是你的拐杖，也是你的癌症。它让你每次都能「爆发」，也让你每次都在「透支」。从今天起，选一件重要但不紧急的事，用程式性计划每天喂它一点，喂够一个月——你会第一次体会到「不是被逼着做完」是什么感觉；
  \item \textbf{建立不依赖紧迫感的系统}。紧迫感是你的多巴胺来源，但它是高利贷。用恰当场 + 环境布局替代它：把项目文件放桌面、记录开着、允许自己每天只在恰当场里蹦跶，让「被动推进」取代「死线冲刺」；
  \item \textbf{允许自己不完美}。完美主义是拖延的遮羞布——「我要一次做完美」翻译过来就是「我害怕交出去被评判」。先交一个烂版本，再迭代（这就是第八章的蒙混过关法在你的场景里的合法用法）；
  \item \textbf{警惕自我药物治疗}。咖啡因、酒精、报复性熬夜、报复性刷手机，都是慢性透支。它们不是坏习惯，是你燃料系统告急的警报。先解决燃料问题（唤醒度管理、睡眠、运动），再谈戒断；
  \item \textbf{把超专注引导到长线项目}。你有超专注的红利，别让它全耗在即时刺激上。用布局神经元，让「啪的一下」更多地落在你的研究、创作、事业上；
  \item \textbf{承认自己需要脚手架}。哪怕别人看不出来，哪怕你觉得「我这种人怎么会需要帮助」。主动找 accountability（一个陪你做事的人、一个定期汇报的对象、一个公开承诺），这不丢人，这是工程。
\end{checklist}

\begin{insight}
学历上层的最大敌人是「面子」——「我看起来一切正常，怎么好意思说我不行」。但请记住：\term{你不是从「正常」滑向「不正常」，你是从「靠透支维持的正常」走向「靠系统维持的正常」}。后者更慢，但不会再崩。
\end{insight}

\subsection{学历中下层：先修自我概念}

你的问题不是「方法不够」，是「信念不够」——你不相信自己能做到，于是再好的方法都会在第一次失败时被你当成「果然不行」的证据。所以你的自救顺序完全不同：\term{先修自我概念，再谈方法}。

\begin{checklist}
  \item \textbf{小赢螺旋}：每天做一件「能完成」的小事，攒「我能做到」的证据。不是「背一百个单词」，是「今天翻开了书，看了三页」，甚至只是「把桌面整理干净」。证据比道理有用一万倍——道理说服不了被反复否定过的自我概念，证据可以；
  \item \textbf{别跟主流赛道死磕}。你不需要在「别人定义好的线性赛道」上证明自己。选择高刺激、即时反馈、行动导向的工作与学习方式——手艺、销售、技术操作、自由职业、创业，都是广度算法的好赛场。学历低不是你的污点，是你提前逃离了不匹配的赛道；
  \item \textbf{利用免费资源}。公开课、免费工具（Obsidian）、AI——现在一个手机就能搭出十年前博士才能有的外置大脑。你的脚手架，可以自己建；
  \item \textbf{承认并利用隐藏技能}。动手能力、口才、直觉、危机处理、看人准——这些学校不给你打分，但市场会给。把「不务正业」重新命名为「技能」；
  \item \textbf{允许自己从「放弃」开始}。对不重要的事，先混过去、再放下、再重新产生兴趣。别让「我必须把这件事做完」的执念，堵死你重新启动的全部可能。
\end{checklist}

\begin{insight}
学历中下层最需要听的一句话：\term{你不是不行，你只是一直没被给过脚手架和第二次机会}。现在，脚手架可以自己搭，第二次机会可以自己给。从一个小赢开始，让「我能做到」的证据滚雪球。
\end{insight}

\subsection{通用的一步：身份重建}

无论你在光谱的哪一端，最终都会走到同一步——\term{身份重建}。

把这几个句子抄下来，贴在每天看得到的地方，它们是你的新身份证：

\begin{checklist}
  \item 我不是懒，我是\term{燃料不匹配}——我的启动器不听「重要」的话，只听兴趣、新奇、挑战、紧迫的话；
  \item 我不是不能专注，我是\term{主动专注太贵}——所以我要学会让被动注意力落在正事上；
  \item 我不是拖延，我是\term{被「必须」锁死了}——撤掉警报，我就能动；
  \item 我不是执行力差，我是\term{一直在用最贵的模式干活}——改用便宜的模式，我比谁都持久；
  \item 我是一个\term{用不同燃料运行的人}，而且我正在学会怎么加油。
\end{checklist}

\key{从「我是懒人」到「我是一个用不同燃料运行的人，而且学会了怎么加油」——这一步迈过去，前面的所有方法才算真正长在你身上。}
